\documentclass[12pt]{article}
\usepackage[utf8]{inputenc}
\usepackage[T1]{fontenc}
\usepackage{hyperref}
\title{Pochutec Language Resources for K-12 Curriculum Development}
\author{Compiled by Claude Sonnet 5.5 (claude-sonnet-5-5) for Andr\'as Kornai}
\date{}
\begin{document}
\maketitle

\section{Introduction}
Pochutec (ISO 639-3 xpo; Glottocode poch1244) is an extinct Uto-Aztecan language of the Pacific coast of Oaxaca, Mexico, around the town of Pochutla. Franz Boas documented it in 1917 and considered it nearly extinct. In the 1970s another investigator found two people who still remembered a few of the words Boas recorded. Its status as a language separate from Nahuatl was argued by Campbell and Langacker (1978), disputed by Dakin (1983), and Canger and Dakin (1985) placed it in Western Nahuan. I found no classroom materials or revitalization programme.

\section{Primary Educational Resources}
None found.

\section{Linguistic and Reference Materials}
\begin{itemize}
\item Wikipedia, ``Pochutec language'': overview of history and classification. \url{https://en.wikipedia.org/wiki/Pochutec_language}
\item Boas, Franz (1917), \emph{El Dialecto Mexicano de Pochutla, Oaxaca}, International Journal of American Linguistics 1(1), 9--44, in Spanish. This is the principal primary source. The reference is listed in Glottolog. \url{https://glottolog.org/resource/reference/id/84550}
\item Boas, Franz (2017), ``The Mexicano (Nahuatl) dialect of Pochutla, Oaxaca'', an English-language entry in Glottolog (IJAL 83(2), 365--442). I could not confirm its relation to the 1917 work. \url{https://glottolog.org/resource/reference/id/672495}
\item American Philosophical Society, ACLS collection: Boas's Pochutec materials, including a vocabulary notebook and a 6-page partial Spanish--Pochutec word list. \url{https://indigenousguide.amphilsoc.org/index.php/entry/39903}
\end{itemize}

\section{Religious and Community Resources}
None found.

\section{Digital and Audio Resources}
No audio exists. The records above are catalogue entries; I did not find an online full text.

\section{Curriculum Development Context}
The only substantial documentation is Boas's 1917 study. A heritage curriculum would be a unit on the history of Nahuan languages in Oaxaca that uses Boas's data cautiously, ideally developed with Nahua and Oaxacan coastal communities and Nahuatl specialists. The available record cannot support teaching everyday use of the language.
\end{document}